\chapter{Quantum Machine Learning}
\label{ch:lec20}

\begin{objectives}
\begin{itemize}[leftmargin=1.4em,nosep]
  \item Load classical data into quantum states and feature maps.
  \item Define and use quantum kernels.
  \item Train variational classifiers / quantum neural networks.
  \item Assess barren plateaus and the practical caveats of QML.
\end{itemize}
\end{objectives}

\todo{Write the introductory paragraph orienting the lecture.}

\section{Data loading and feature maps}
\label{sec:lec20-1}
\todo{Draft this section.}

\section{Quantum kernels}
\label{sec:lec20-2}
\todo{Draft this section.}

\section{Variational classifiers and QNNs}
\label{sec:lec20-3}
\todo{Draft this section.}

\section{Linear-algebra-based QML (HHL, PCA)}
\label{sec:lec20-4}
\todo{Draft this section.}

\section{Barren plateaus and caveats}
\label{sec:lec20-5}
\todo{Draft this section.}

\exercisesection
\begin{exercise}
\todo{Add exercise.}
\end{exercise}
\begin{exercise}
\todo{Add exercise.}
\end{exercise}
\begin{exercise}
\todo{Add exercise.}
\end{exercise}

\readingsection
\begin{itemize}[leftmargin=1.4em]
  \item M.~Schuld, N.~Killoran, \emph{Quantum Machine Learning in Feature Hilbert Spaces}, PRL 122 (2019).
  \item J.~Biamonte et al., \emph{Quantum Machine Learning}, Nature 549 (2017).
\end{itemize}
